\documentclass[10pt,a4paper]{amsart}
\usepackage{amsmath,amssymb,amsthm,graphicx}
% Retained minimal preamble for the Stacks Project chapter coherent.tex (Stacks chapter, pre-delegation retained fragment for stacks-coherent).
% Verbatim excerpt of the chapter preamble (sources/stacks-spaces-descent/preamble.tex): the reference/slogan/history
% environment block (native lines 7-12), Xy-pic setup (native lines 14-19) and the native theorem declarations with
% the mathematical macros (native lines 154-217). Copied mechanically; nothing added, removed or redefined. The
% wrapper class and the amsmath/amssymb/amsthm/graphicx base packages are supplied by the assembler, so the
% class/xr-hyper/hyperref/externaldocument part of the native preamble is omitted.
% For dealing with references we use the comment environment
\usepackage{verbatim}
\newenvironment{reference}{\comment}{\endcomment}
%\newenvironment{reference}{}{}
\newenvironment{slogan}{\comment}{\endcomment}
\newenvironment{history}{\comment}{\endcomment}
% For commutative diagrams we use Xy-pic
\usepackage[all]{xy}

% We use 2cell for 2-commutative diagrams.
\xyoption{2cell}
\UseAllTwocells
% Theorem environments.
%
\theoremstyle{plain}
\newtheorem{theorem}[subsection]{Theorem}
\newtheorem{proposition}[subsection]{Proposition}
\newtheorem{lemma}[subsection]{Lemma}

\theoremstyle{definition}
\newtheorem{definition}[subsection]{Definition}
\newtheorem{example}[subsection]{Example}
\newtheorem{exercise}[subsection]{Exercise}
\newtheorem{situation}[subsection]{Situation}

\theoremstyle{remark}
\newtheorem{remark}[subsection]{Remark}
\newtheorem{remarks}[subsection]{Remarks}

\numberwithin{equation}{subsection}

% Macros
%
\def\lim{\mathop{\mathrm{lim}}\nolimits}
\def\colim{\mathop{\mathrm{colim}}\nolimits}
\def\Spec{\mathop{\mathrm{Spec}}}
\def\Hom{\mathop{\mathrm{Hom}}\nolimits}
\def\Ext{\mathop{\mathrm{Ext}}\nolimits}
\def\SheafHom{\mathop{\mathcal{H}\!\mathit{om}}\nolimits}
\def\SheafExt{\mathop{\mathcal{E}\!\mathit{xt}}\nolimits}
\def\Sch{\mathit{Sch}}
\def\Mor{\mathop{\mathrm{Mor}}\nolimits}
\def\Ob{\mathop{\mathrm{Ob}}\nolimits}
\def\Sh{\mathop{\mathit{Sh}}\nolimits}
\def\NL{\mathop{N\!L}\nolimits}
\def\CH{\mathop{\mathrm{CH}}\nolimits}
\def\proetale{{pro\text{-}\acute{e}tale}}
\def\etale{{\acute{e}tale}}
\def\QCoh{\mathit{QCoh}}
\def\Ker{\mathop{\mathrm{Ker}}}
\def\Im{\mathop{\mathrm{Im}}}
\def\Coker{\mathop{\mathrm{Coker}}}
\def\Coim{\mathop{\mathrm{Coim}}}

% Boxtimes
%
\DeclareMathSymbol{\boxtimes}{\mathbin}{AMSa}{"02}

%
% Macros for moduli stacks/spaces
%
\def\QCohstack{\mathcal{QC}\!\mathit{oh}}
\def\Cohstack{\mathcal{C}\!\mathit{oh}}
\def\Spacesstack{\mathcal{S}\!\mathit{paces}}
\def\Quotfunctor{\mathrm{Quot}}
\def\Hilbfunctor{\mathrm{Hilb}}
\def\Curvesstack{\mathcal{C}\!\mathit{urves}}
\def\Polarizedstack{\mathcal{P}\!\mathit{olarized}}
\def\Complexesstack{\mathcal{C}\!\mathit{omplexes}}
% \Pic is the operator that assigns to X its picard group, usage \Pic(X)
% \Picardstack_{X/B} denotes the Picard stack of X over B
% \Picardfunctor_{X/B} denotes the Picard functor of X over B
\def\Pic{\mathop{\mathrm{Pic}}\nolimits}
\def\Picardstack{\mathcal{P}\!\mathit{ic}}
\def\Picardfunctor{\mathrm{Pic}}
\def\Deformationcategory{\mathcal{D}\!\mathit{ef}}
\usepackage{cleveref}
\crefname{originallemma}{original lemma}{original lemmas}
\Crefname{originallemma}{Original lemma}{Original lemmas}
\makeatletter
\newlabel{lemma-higher-direct-images-zero-finite-fibre}{{02OE}{source}{}{sourcepointer-1}{}}
\newlabel{lemma-higher-direct-images-zero-finite-fibre@cref}{{[originallemma][0][]02OE}{[1][0][]source}{}{}{}}
\makeatother
\crefname{originallemma}{original lemma}{original lemmas}
\Crefname{originallemma}{Original lemma}{Original lemmas}
\makeatletter
\newlabel{lemma-quasi-coherent-affine-cohomology-zero}{{01XB}{source}{}{sourcepointer-2}{}}
\newlabel{lemma-quasi-coherent-affine-cohomology-zero@cref}{{[originallemma][0][]01XB}{[1][0][]source}{}{}{}}
\makeatother
\crefname{originallemma}{original lemma}{original lemmas}
\Crefname{originallemma}{Original lemma}{Original lemmas}
\makeatletter
\newlabel{lemma-quasi-separated-h1-zero-covering}{{01XG}{source}{}{sourcepointer-3}{}}
\newlabel{lemma-quasi-separated-h1-zero-covering@cref}{{[originallemma][0][]01XG}{[1][0][]source}{}{}{}}
\makeatother
\crefname{originallemma}{original lemma}{original lemmas}
\Crefname{originallemma}{Original lemma}{Original lemmas}
\makeatletter
\newlabel{more-morphisms-lemma-characterize-finite}{{02LS}{source}{}{sourcepointer-4}{}}
\newlabel{more-morphisms-lemma-characterize-finite@cref}{{[originallemma][0][]02LS}{[1][0][]source}{}{}{}}
\makeatother
\crefname{originallemma}{original lemma}{original lemmas}
\Crefname{originallemma}{Original lemma}{Original lemmas}
\makeatletter
\newlabel{morphisms-lemma-finite-integral}{{01WJ}{source}{}{sourcepointer-5}{}}
\newlabel{morphisms-lemma-finite-integral@cref}{{[originallemma][0][]01WJ}{[1][0][]source}{}{}{}}
\makeatother
\crefname{originallemma}{original lemma}{original lemmas}
\Crefname{originallemma}{Original lemma}{Original lemmas}
\makeatletter
\newlabel{morphisms-lemma-finite-proper}{{01WN}{source}{}{sourcepointer-6}{}}
\newlabel{morphisms-lemma-finite-proper@cref}{{[originallemma][0][]01WN}{[1][0][]source}{}{}{}}
\makeatother
\crefname{originallemma}{original lemma}{original lemmas}
\Crefname{originallemma}{Original lemma}{Original lemmas}
\makeatletter
\newlabel{morphisms-lemma-finite-quasi-finite}{{02NU}{source}{}{sourcepointer-7}{}}
\newlabel{morphisms-lemma-finite-quasi-finite@cref}{{[originallemma][0][]02NU}{[1][0][]source}{}{}{}}
\makeatother
\crefname{originallemma}{original lemma}{original lemmas}
\Crefname{originallemma}{Original lemma}{Original lemmas}
\makeatletter
\newlabel{morphisms-lemma-integral-universally-closed}{{01WM}{source}{}{sourcepointer-8}{}}
\newlabel{morphisms-lemma-integral-universally-closed@cref}{{[originallemma][0][]01WM}{[1][0][]source}{}{}{}}
\makeatother
\crefname{originallemma}{original lemma}{original lemmas}
\Crefname{originallemma}{Original lemma}{Original lemmas}
\makeatletter
\newlabel{morphisms-lemma-quasi-finite}{{02NH}{source}{}{sourcepointer-9}{}}
\newlabel{morphisms-lemma-quasi-finite@cref}{{[originallemma][0][]02NH}{[1][0][]source}{}{}{}}
\makeatother
\crefname{originallemma}{original lemma}{original lemmas}
\Crefname{originallemma}{Original lemma}{Original lemmas}
\makeatletter
\newlabel{schemes-lemma-push-forward-quasi-coherent}{{01LC}{source}{}{sourcepointer-10}{}}
\newlabel{schemes-lemma-push-forward-quasi-coherent@cref}{{[originallemma][0][]01LC}{[1][0][]source}{}{}{}}
\makeatother
\title{Finiteness via properness and finite fibres over a locally Noetherian base: for f : X -> S with S locally Noetherian, f is finite if and only if f is proper with finite fibres (retained with the author's parenthetical pointing to the more general More on Morphisms twin).}
\author{The Stacks Project Authors}
\date{stacks-project:coherent.tex:local-snapshot-20261012; GNU FDL 1.2 or later}
\begin{document}
\maketitle
\noindent\textit{Editorial note.} Statement span is lines 6097-6108 of sources/stacks-coherent/pinned-original-source.tex; proof span is lines 6110-6149, the complete author proof (38 body lines).\par\medskip
\noindent\textit{Editorial note.} Navigation stubs: 10 cited-but-not-included labels are supplied as non-load-bearing pointers with their pinned tags where the snapshot tag index has them.\par\medskip
\noindent\textit{Editorial note.} extra\_packages is empty on purpose: the retained preamble already loads xy with options and a second declaration clashes in the assembler wrapper.\par\medskip
\noindent\textit{Editorial note.} Text is retained verbatim, including the author's own typos and placeholders.\par\medskip
\noindent\textit{Editorial note.} Pool stacks-coherent backfill wave 33c. precard\_gate RESULT=OK (probe L9999).\par\medskip
\noindent\textit{Editorial note.} Cross-chapter same-label warning (NOT an occupancy conflict, Chief ruling 2026-10-13 evening): label 'lemma-characterize-finite' also exists in stacks-spaces-more-morphisms — a different proposition in a different chapter file.\par\medskip
\noindent\textit{Allowed standard prerequisites (editorial).}\begin{itemize}
\item Explicit prerequisite boundary. See why-field; every pointer named here occurs verbatim in the retained statement/proof text.
\end{itemize}\medskip
\ifdefined\hypertarget\hypertarget{sourcepointer-1}{}\fi
\noindent\textit{Source reference.} Original lemma 02OE: sources/stacks-coherent/pinned-original-source.tex, source lines 6013--6021. Cited by the retained statement or proof but not part of the collected spans; supplied only as a navigation pointer. The argument does not depend on its content.\par\medskip
\ifdefined\hypertarget\hypertarget{sourcepointer-2}{}\fi
\noindent\textit{Source reference.} Original lemma 01XB: sources/stacks-coherent/pinned-original-source.tex, source lines 363--367. Cited by the retained statement or proof but not part of the collected spans; supplied only as a navigation pointer. The argument does not depend on its content.\par\medskip
\ifdefined\hypertarget\hypertarget{sourcepointer-3}{}\fi
\noindent\textit{Source reference.} Original lemma 01XG: sources/stacks-coherent/pinned-original-source.tex, source lines 608--611. Cited by the retained statement or proof but not part of the collected spans; supplied only as a navigation pointer. The argument does not depend on its content.\par\medskip
\ifdefined\hypertarget\hypertarget{sourcepointer-4}{}\fi
\noindent\textit{Source reference.} Original lemma 02LS: sources/stacks-coherent/pinned-original-source.tex, source lines 6149--6149. Cited by the retained statement or proof but not part of the collected spans; supplied only as a navigation pointer. The argument does not depend on its content.\par\medskip
\ifdefined\hypertarget\hypertarget{sourcepointer-5}{}\fi
\noindent\textit{Source reference.} Original lemma 01WJ: sources/stacks-coherent/pinned-original-source.tex, source lines 6149--6149. Cited by the retained statement or proof but not part of the collected spans; supplied only as a navigation pointer. The argument does not depend on its content.\par\medskip
\ifdefined\hypertarget\hypertarget{sourcepointer-6}{}\fi
\noindent\textit{Source reference.} Original lemma 01WN: sources/stacks-coherent/pinned-original-source.tex, source lines 6149--6149. Cited by the retained statement or proof but not part of the collected spans; supplied only as a navigation pointer. The argument does not depend on its content.\par\medskip
\ifdefined\hypertarget\hypertarget{sourcepointer-7}{}\fi
\noindent\textit{Source reference.} Original lemma 02NU: sources/stacks-coherent/pinned-original-source.tex, source lines 6149--6149. Cited by the retained statement or proof but not part of the collected spans; supplied only as a navigation pointer. The argument does not depend on its content.\par\medskip
\ifdefined\hypertarget\hypertarget{sourcepointer-8}{}\fi
\noindent\textit{Source reference.} Original lemma 01WM: sources/stacks-coherent/pinned-original-source.tex, source lines 6149--6149. Cited by the retained statement or proof but not part of the collected spans; supplied only as a navigation pointer. The argument does not depend on its content.\par\medskip
\ifdefined\hypertarget\hypertarget{sourcepointer-9}{}\fi
\noindent\textit{Source reference.} Original lemma 02NH: sources/stacks-coherent/pinned-original-source.tex, source lines 6149--6149. Cited by the retained statement or proof but not part of the collected spans; supplied only as a navigation pointer. The argument does not depend on its content.\par\medskip
\ifdefined\hypertarget\hypertarget{sourcepointer-10}{}\fi
\noindent\textit{Source reference.} Original lemma 01LC: sources/stacks-coherent/pinned-original-source.tex, source lines 6149--6149. Cited by the retained statement or proof but not part of the collected spans; supplied only as a navigation pointer. The argument does not depend on its content.\par\medskip
\begin{lemma}
\label{lemma-characterize-finite}
(For a more general version see
More on Morphisms, Lemma \ref{more-morphisms-lemma-characterize-finite}.)
Let $f : X \to S$ be a morphism of schemes.
Assume $S$ is locally Noetherian.
The following are equivalent
\begin{enumerate}
\item $f$ is finite, and
\item $f$ is proper with finite fibres.
\end{enumerate}
\end{lemma}

\begin{proof}
A finite morphism is proper according to
Morphisms, Lemma \ref{morphisms-lemma-finite-proper}.
A finite morphism is quasi-finite according to
Morphisms, Lemma \ref{morphisms-lemma-finite-quasi-finite}.
A quasi-finite morphism has finite fibres, see
Morphisms, Lemma \ref{morphisms-lemma-quasi-finite}.
Hence a finite morphism is proper and has finite fibres.

\medskip\noindent
Assume $f$ is proper with finite fibres.
We want to show $f$ is finite.
In fact it suffices to prove $f$ is affine.
Namely, if $f$ is affine, then it follows that
$f$ is integral by
Morphisms, Lemma \ref{morphisms-lemma-integral-universally-closed}
whereupon it follows from
Morphisms, Lemma \ref{morphisms-lemma-finite-integral}
that $f$ is finite.

\medskip\noindent
To show that $f$ is affine we may assume that $S$ is affine, and our
goal is to show that $X$ is affine too.
Since $f$ is proper we see that $X$ is separated and quasi-compact.
Hence we may use the criterion of
Lemma \ref{lemma-quasi-separated-h1-zero-covering} to prove that $X$
is affine. To see this let $\mathcal{I} \subset \mathcal{O}_X$
be a finite type ideal sheaf. In particular $\mathcal{I}$ is
a coherent sheaf on $X$. By
Lemma \ref{lemma-higher-direct-images-zero-finite-fibre} we conclude that
$R^1f_*\mathcal{I}_s = 0$ for all $s \in S$.
In other words, $R^1f_*\mathcal{I} = 0$. Hence we see from
the Leray Spectral Sequence for $f$ that
$H^1(X , \mathcal{I}) = H^1(S, f_*\mathcal{I})$.
Since $S$ is affine, and $f_*\mathcal{I}$ is quasi-coherent
(Schemes, Lemma \ref{schemes-lemma-push-forward-quasi-coherent})
we conclude $H^1(S, f_*\mathcal{I}) = 0$
from Lemma \ref{lemma-quasi-coherent-affine-cohomology-zero}
as desired. Hence $H^1(X, \mathcal{I}) = 0$ as desired.
\end{proof}

\end{document}
